% ============================================================================
% REFERENCIAS BIBLIOGRAFICAS
% ============================================================================

\section*{REFER\^ENCIAS BIBLIOGR\'AFICAS}
\addcontentsline{toc}{section}{Refer\^encias Bibliogr\'aficas}

\begin{description}[style=nextline, leftmargin=0pt, labelindent=0pt]

\item[ABNT, 2009a.] ASSOCIA\c{C}\~AO BRASILEIRA DE NORMAS T\'ECNICAS. \textbf{NBR ISO 14040:2009} --- Gest\~ao ambiental --- Avalia\c{c}\~ao do ciclo de vida --- Princ\'ipios e estrutura. Rio de Janeiro, 2009.

\item[ABNT, 2009b.] ASSOCIA\c{C}\~AO BRASILEIRA DE NORMAS T\'ECNICAS. \textbf{NBR ISO 14044:2009} --- Gest\~ao ambiental --- Avalia\c{c}\~ao do ciclo de vida --- Requisitos e orienta\c{c}\~oes. Rio de Janeiro, 2009.

\item[Ecoinvent, 2024.] ECOINVENT CENTRE. \textbf{Ecoinvent database v3}. Swiss Centre for Life Cycle Inventories, 2024. Dispon\'ivel em: \url{https://ecoinvent.org/}

\item[GreenDelta, 2024.] GREENDELTA GmbH. \textbf{OpenLCA} --- Open source life cycle assessment software. Berlin, 2024. Dispon\'ivel em: \url{https://www.openlca.org/}

\item[Giusti, 2025.] GIUSTI, Gabriela. \textbf{Development of characterization factors for health effects of particulate matter in Brazil}. 2025. Tese (Doutorado em Planejamento e Uso de Recursos Renov\'aveis) --- Universidade Federal de S\~ao Carlos, Sorocaba, 2025. Dispon\'ivel em: \url{https://repositorio.ufscar.br/handle/20.500.14289/22123}

\item[Hauschild \& Huijbregts, 2015.] HAUSCHILD, Michael Z.; HUIJBREGTS, Mark A. J. (Ed.). \textbf{Life Cycle Impact Assessment}. Springer, 2015. (LCA Compendium --- The Complete World of Life Cycle Assessment).

\item[Huijbregts et al., 2017.] HUIJBREGTS, Mark A. J. et al. ReCiPe2016: a harmonised life cycle impact assessment method at midpoint and endpoint level. \textbf{International Journal of Life Cycle Assessment}, v. 22, p. 138--147, 2017.

\item[Rosenbaum et al., 2008.] ROSENBAUM, Ralph K. et al. USEtox --- the UNEP-SETAC toxicity model: recommended characterisation factors for human toxicity and freshwater ecotoxicity in life cycle impact assessment. \textbf{International Journal of Life Cycle Assessment}, v. 13, p. 532--546, 2008.

\item[Verones et al., 2020.] VERONES, Francesca et al. \textbf{LC-Impact Version 1.0 --- A spatially differentiated life cycle impact assessment approach}. 2020. Dispon\'ivel em: \url{https://lc-impact.eu/}

\end{description}
